\documentclass[../../../include/open-logic-section]{subfiles}

\begin{document}

\olfileid{sth}{ordinals}{wo}
\olsection{સુક્રમો}

મૂળભૂત ખ્યાલ આ છે:

\begin{defn}
સંબંધ $<$ એ $A$ને \emph{સુક્રમિત} કરે જો અને ફક્ત જો નીચેની
બે શરતો પૂરી થાય:
\begin{enumerate}
	\item $<$ તુલનાયુક્ત છે, એટલે કે બધા $a, b \in A$ માટે
	કાં $a < b$ અથવા $a = b$ અથવા $b < a$;
	\item $A$ના દરેક અરિક્ત ઉપગણને $<$ મુજબ ન્યૂનતમ !!{element}
	હોય, એટલે કે જો $\emptyset \neq X \subseteq A$ તો
	$(\exists m \in
	X)(\forall z \in X)z \nless m$
\end{enumerate}
\end{defn}

હમણાં વિચારેલા ત્રણેય દાખલા ખરેખર સુક્રમિત કરનારા સંબંધો છે
તે સહેલાઈથી જોઈ શકાય છે.

\begin{prob}
\Olref[sth][ordinals][idea]{sec}માં પ્રાકૃતિક સંખ્યાઓ પરના
ત્રણ ક્રમોના દાખલા આપ્યા હતા. દરેક સુક્રમ છે તે તપાસો.
\end{prob}

સુક્રમ વિશેનાં કેટલાંક પ્રાથમિક પણ અત્યંત મહત્ત્વનાં અવલોકનો આ છે.

\begin{prop}\ollabel{wo:strictorder}
જો $<$ એ $A$ને સુક્રમિત કરતો હોય, તો $A$ના દરેક અરિક્ત
ઉપગણને અનન્ય $<$ મુજબ લઘુતમ સભ્ય હોય, અને $<$ અસ્વવાચક,
અસંમિત તથા પરંપરિત હોય.
\end{prop}

\begin{proof}
જો $X$ એ $A$નો અરિક્ત ઉપગણ હોય, તો તેને $<$ મુજબ
ન્યૂનતમ !!{element} $m$ હોય, એટલે કે
$(\forall z \in X)z \nless m$.
$<$ તુલનાયુક્ત હોવાથી $(\forall z \in X)m \leq z$.
તેથી $m$ એ $X$નો $<$ મુજબ લઘુતમ !!{element} છે.

અસ્વવાચકતા માટે કોઈ $a \in A$ લો; $\{a\}$નો $<$ મુજબ
લઘુતમ !!{element} $a$ છે, તેથી $a \nless a$.
પરંપરિતતા માટે, જો $a < b < c$, તો
$\{a, b, c\}$ને $<$ મુજબ લઘુતમ !!{element} હોવાથી
$a < c$. અસંમિતતા અસ્વવાચકતા અને પરંપરિતતાથી
મળે છે.
\end{proof}

\begin{prop}\ollabel{propwoinduction}
જો $<$ એ $A$ને સુક્રમિત કરતો હોય, તો કોઈ પણ સૂત્ર
$\phi(x)$ માટે:
\[
	\text{જો }(\forall a \in A)((\forall b < a)\phi(b) \lif
		\phi(a))\text{, તો }(\forall a \in A)\phi(a).
\]
\end{prop}

\begin{proof}
આપણે પ્રતિવિધાન સાબિત કરીશું. ધારો કે $\lnot(\forall a \in
A)\phi(a)$, એટલે કે $X = \Setabs{x \in A}{\lnot\phi(x)} \neq
\emptyset$. ત્યારે $X$ને $<$ મુજબ ન્યૂનતમ !!{element} $a$
હોય. તેથી $(\forall
b < a)\phi(b)$, પરંતુ $\lnot \phi(a)$.
\end{proof}\noindent આ છેલ્લા ગુણધર્મથી પ્રાકૃતિક સંખ્યાઓ પરના
પ્રબળ ગાણિતિક અનુમાનનો સિદ્ધાંત યાદ આવવો જોઈએ, એટલે કે:
જો $(\forall n \in \omega)((\forall m < n)\phi(m)
\lif \phi(n))$ હોય, તો $(\forall n \in \omega)\phi(n)$.
આ ગુણધર્મ સુક્રમને ખૂબ \emph{મજબૂત} ખ્યાલ બનાવે છે.\footnote{
યાદ રાખો: સ્પષ્ટપણે વિરુદ્ધ ન કહ્યું હોય તો બધા સૂત્રોમાં
પ્રાચલો હોઈ શકે છે.}

\end{document}
